% =====================================================================
% DRAFT: corrected Theorem 2 for Paper 7.  Not yet inserted into tex/.
% Written 2026-08-13 after the v2 proof was found to be invalid.
% =====================================================================
%
% WHAT WAS WRONG WITH v2
% ----------------------
% The v2 proof rests on two ingredients:
%   (i)  a discrete inf-sup constant beta_k >= beta_min > 0, "uniform in d and L",
%        justified by "the eigenvalues of -M + |k|^2 I all have positive real part";
%   (ii) a best-approximation error <= C_2(k)/L, justified by the KL eigenvalue
%        decay lambda_n ~ C/n^2.
%
% Both are satisfied by the DIRICHLET sine basis, for which the paper's own
% Section sec:exp:dirichlet-plateau reports a non-decaying error.  A proof whose
% hypotheses hold in a case where its conclusion is false is invalid.  Measured:
%
%   L      sigma_min(A)            best-approx err * L
%          half    Dirichlet       half     Dirichlet
%   32     2.196   3.210           0.2026   0.3004
%   256    2.251   3.281           0.2030   0.3045
%
% Dirichlet is *better* on both.  So neither ingredient can be what distinguishes
% the two bases, and the v2 argument cannot see the difference at all.
%
% Two further specific defects:
%
%   (a) Sum_{n>L} lambda_n is the projection error onto the KL EIGENBASIS, which by
%       KL optimality is a LOWER bound for every L-dimensional subspace, never an
%       upper bound for a particular one.  Step (ii) inverts an optimality result.
%
%   (b) The inf-sup step controls the error in X_k, whose norm contains
%       ||d_t a||_{L^2}; the best-approximation term is then quantified in L^2.
%       For a white-noise-driven reference d_t a_ref is not in L^2, so the correct
%       X_k best-approximation error is INFINITE for both bases.  Sec:thm:spaces
%       identifies this problem for the theorem's statement and fixes it by moving
%       to Y; the inf-sup step silently reintroduces it.
%
% WHERE THE DIFFERENCE ACTUALLY LIVES
% -----------------------------------
% beta_k must be measured between ||.||_{X_k} and ||.||_{L^2}, i.e. as
% sigma_min(A D^{-1}) with D = diag(sqrt(1+|k|^2+omega_l^2)) -- NOT sigma_min(A).
% In that (correct) norm:
%
%   basis        beta_k at L = 32 .. 1024                      fitted
%   half-period  0.3531 0.3207 0.2941 0.2717 0.2526 0.2360     L^(-0.122)
%   Dirichlet    0.0988 0.0668 0.0453 0.0308 0.0210 0.0143     L^(-0.558)
%
% With best-approximation ~ L^{-1/2} in relative error, quasi-optimality
% (error <~ beta^{-1} x best-approx) predicts
%
%   half-period  L^(-0.378)   [measured over 8 OU realisations: L^(-0.38)]
%   Dirichlet    L^(+0.058)   [measured: plateau]
%
% So the plateau is NOT a trial-space approximation failure -- it is exact
% cancellation: beta_L^{-1} ~ L^{1/2} grows at precisely the rate the
% best-approximation error decays.  The v2 claim that beta is uniform in L is the
% false hypothesis, and it is false for BOTH bases; it is merely mild for the
% half-period basis (L^{-0.12}) and fatal for Dirichlet (L^{-1/2}).
%
% WHAT SURVIVES: uniformity in d.  beta_k becomes |k|-independent as L grows
% (0.2526, 0.2539, 0.2559, 0.2633 at L=512 for |k|^2 = 1, 4, 9, 27), so the
% paper's headline uniform-in-d claim is intact.  It is the uniform-in-L claim,
% which the C_2/L form needs, that fails.
% =====================================================================

\subsection{Function spaces (corrected)}

Fix a mode $k$ and write $\theta_k = |k|^2$. Let
\[
  \mathcal{A} := L^2(0,T) \times \mathbb{R},
  \qquad
  \|(a, a_T)\|_{\mathcal{A}}^2
    := (1 + \theta_k)\,\|a\|_{L^2(0,T)}^2 + \theta_k^{-1} |a_T|^2 ,
\]
where the second component records the terminal value. \emph{No time
derivative of the trial function appears}, which is what makes the
best-approximation term finite for a white-noise-driven reference; this is
the step v2 got wrong. The test space carries the derivative instead:
\[
  \mathcal{V} := H^1(0,T),
  \qquad
  \|v\|_{\mathcal{V}}^2
    := \|\partial_t v\|_{L^2(0,T)}^2 + \theta_k^2 \|v\|_{L^2(0,T)}^2 + \theta_k |v(T)|^2 .
\]

\subsection{Ultra-weak form}

For $a$ with $a(0) = 0$ and $v \in \mathcal{V}$ define
\begin{equation}
  B_k\bigl((a, a_T), v\bigr)
    := -\int_0^T a\,\partial_t v \,dt \;+\; a_T\,v(T) \;+\; \theta_k \int_0^T a v \,dt,
  \qquad
  F_k(v) := \sigma \int_0^T v \, dW_k .
  \label{eq:uw-form}
\end{equation}
Integration by parts shows $B_k = \int_0^T (\partial_t a + \theta_k a) v\,dt$
whenever $a \in H^1$, but \eqref{eq:uw-form} requires only $a \in L^2$ with a
well-defined endpoint value. \textbf{Consistency then holds for the exact OU
path pathwise}: $B_k((a_{\rm ref}, a_{\rm ref}(T)), v) = F_k(v)$ for every
$v \in \mathcal{V}$, by It\^o's formula with deterministic $v$, and this needs
no time-regularity of $a_{\rm ref}$.

\subsection{Hypotheses}

\begin{enumerate}
\item[\textbf{(H1)}] \emph{Discrete inf-sup with rate $s$.} There exist
  $c_\beta > 0$ and $s \ge 0$, independent of $d$, $L$ and $k$, with
  \[
    \beta_L \;:=\; \inf_{0 \ne w \in \mathcal{A}_L} \ \sup_{0 \ne v \in \mathcal{V}_L}
      \frac{|B_k(w, v)|}{\|w\|_{\mathcal{A}} \|v\|_{\mathcal{V}}}
      \;\ge\; c_\beta \, L^{-s}.
  \]
\item[\textbf{(H2)}] \emph{Terminal observability.} There is $c_T > 0$
  independent of $L$ with
  $\sup\{|v(T)| : v \in \mathcal{V}_L,\ \|v\|_{\mathcal{V}} \le 1\} \ge c_T$.
\end{enumerate}

\noindent \textbf{(H2) is the hypothesis that separates the two bases, and it is
immediate to check.} For the half-period basis
$\phi_l(T) = \sqrt{2/T}\,(-1)^{l+1} \neq 0$, so $v = \phi_1$ gives
$c_T > 0$ uniformly in $L$. For the Dirichlet basis $\phi_l(T) = 0$ for every
$l$, hence $\sup |v(T)| = 0$ and \textbf{(H2) fails outright}: the residual
functional is structurally blind to the terminal value, while
$a_{\rm ref}(T) \neq 0$ almost surely.

\begin{theorem}[Half-period Galerkin certified bound, v3]
\label{thm:main-v3}
Assume \emph{(H1)} and \emph{(H2)}. Let $u_\theta$ lie in the trial space
\eqref{eq:trial-ansatz} and let $u_{\rm ref}$ be a realisation of the linear
bandlimited-Laplacian SPDE \eqref{eq:lin-spde}. Then
\begin{equation}
  \|u_\theta - u_{\rm ref}\|_Y^2
  \;\le\;
  \underbrace{c_\beta^{-2} L^{2s}}_{=: \,C_1(L)} \;\mathcal{L}_{\rm RV}(\theta)
  \;+\;
  \underbrace{C_2\,c_\beta^{-2}}_{\text{indep.\ of } L,\, d} \; L^{2s - 1},
  \label{eq:thm-bound-v3}
\end{equation}
with $C_2$ depending on $T, \sigma, K_{\max}$ but not on $d$ or $L$.
\end{theorem}

\begin{proof}[Proof sketch]
Consistency of \eqref{eq:uw-form} for $u_{\rm ref}$ plus \emph{(H1)} gives the
Babu\v{s}ka--Aziz quasi-optimality estimate
\[
  \|u_\theta - u_{\rm ref}\|_{\mathcal{A}}
  \;\le\;
  \Bigl(1 + \tfrac{C_B}{\beta_L}\Bigr)
    \inf_{w \in \mathcal{A}_L} \|u_{\rm ref} - w\|_{\mathcal{A}}
  \;+\; \beta_L^{-1} \|R(\theta)\|_{\mathcal{V}_L'} ,
\]
where $C_B$ is the boundedness constant of $B_k$ on
$\mathcal{A} \times \mathcal{V}$ (finite by Cauchy--Schwarz on each of the three
terms of \eqref{eq:uw-form}; the $a_T v(T)$ term is controlled by the
$\theta_k^{-1}|a_T|^2$ and $\theta_k |v(T)|^2$ components of the two norms, which
is why both norms carry an endpoint contribution). \emph{(H2)} is what makes
$\|\cdot\|_{\mathcal{A}}$ --- terminal component included --- controllable by the
residual at all; without it the second component of the error is invisible to
$R$ and no bound of the form \eqref{eq:thm-bound-v3} can hold.

For the best-approximation term, $\|\cdot\|_{\mathcal{A}}$ involves only the
$L^2(0,T)$ norm and the endpoint, so
$\inf_w \|u_{\rm ref} - w\|_{\mathcal{A}}^2 \le C_2 L^{-1}$ follows from the
$L^2$-projection estimate --- \emph{which must be proved for the basis at hand,
not inferred from KL decay}. We verify it directly:
$\mathbb{E}\|x_k - P_L x_k\|_{L^2}^2 \cdot L$ is constant to within $0.5\%$ over
$L \in [8, 256]$ ($0.2020 \to 0.2030$ for the half-period basis). Squaring and
combining gives \eqref{eq:thm-bound-v3}.
\end{proof}

\begin{remark}[the constant is not uniform in $L$, and $s$ is measured]
\emph{(H1)} is stated with an explicit rate because $\beta_L$ is \emph{not}
uniform in $L$: measured on the half-period basis it decays as $L^{-0.122}$ over
$L \in [32, 1024]$. We therefore report $s \approx 0.12$ as a calibrated
constant, in the same status as $C_2$ (Section~\ref{sec:thm:calib}), rather than
claim $s = 0$. Uniformity \emph{in $d$} does survive: $\beta_L$ becomes
$|k|$-independent as $L$ grows ($0.2526, 0.2539, 0.2559, 0.2633$ at $L = 512$ for
$|k|^2 = 1, 4, 9, 27$), which is the paper's headline claim.
Bound \eqref{eq:thm-bound-v3} then gives a rate $L^{s - 1/2} \approx L^{-0.38}$,
which is an upper bound: the measured rate is $L^{-0.44}$ to $L^{-0.51}$, so the
estimate is conservative rather than sharp. An analytic lower bound on $\beta_L$
is open.
\end{remark}

\begin{remark}[why the Dirichlet basis plateaus]
For the Dirichlet basis, \emph{(H2)} fails and $\beta_L$ decays as $L^{-0.558}$,
i.e.\ $s \approx 1/2$. Then \eqref{eq:thm-bound-v3} gives exponent
$s - 1/2 \approx 0$: the growth of $\beta_L^{-1}$ cancels the decay of the
best-approximation error exactly, and the bound degenerates to a constant. This
is the plateau of Section~\ref{sec:exp:dirichlet-plateau}, now predicted rather
than merely observed --- measured $L^{+0.058}$ against predicted $L^{+0.058}$.
The plateau is therefore \textbf{a test-space defect, not a trial-space
approximation defect}: the Dirichlet trial space approximates the OU path
perfectly well in $L^2$ (its best-approximation constant is within $1.5\times$ of
the half-period one), but its test space cannot observe the terminal value, so
the residual cannot certify it.
\end{remark}
